\documentclass[11pt]{article}
\usepackage{amsmath,amssymb,mathtools}
\usepackage[a4paper,margin=1in]{geometry}

\begin{document}

\section*{Problem 2: Solving the KKT system (0.5 pt)}

Solve the following convex quadratic problem in $\mathbb{R}^3$ with linear equality constraints:
\begin{equation*}
\begin{aligned}
\min_{x\in\mathbb{R}^3}\quad & f(x)
= (-3\ -2\ 7)x
+ \tfrac12 x^\top
\begin{pmatrix}
6 & 0 & 2\\
0 & 4 & 0\\
2 & 0 & 6
\end{pmatrix}
x \\[0.4em]
\text{s.t.}\quad &
\begin{pmatrix}
1 & 1 & 2\\
0 & 1 & 1
\end{pmatrix}x =
\begin{pmatrix}
0\\
2
\end{pmatrix}.
\end{aligned}
\end{equation*}
\noindent
\textit{Hint:} use Gaussian elimination or QR decomposition to solve the resulting KKT system.

\paragraph{Solution 2.}

We write the problem in the standard quadratic form
\begin{equation*}
\min_{x\in\mathbb{R}^3} \ \tfrac12 x^\top Q x + c^\top x
\quad\text{s.t.}\quad Cx = d,
\end{equation*}
with
\begin{equation*}
Q =
\begin{pmatrix}
6 & 0 & 2\\
0 & 4 & 0\\
2 & 0 & 6
\end{pmatrix},\qquad
c =
\begin{pmatrix}
-3\\
-2\\
7
\end{pmatrix},\qquad
C =
\begin{pmatrix}
1 & 1 & 2\\
0 & 1 & 1
\end{pmatrix},\qquad
d =
\begin{pmatrix}
0\\
2
\end{pmatrix}.
\end{equation*}
Introducing Lagrange multipliers $\lambda\in\mathbb{R}^2$ for the equality constraints, the KKT conditions are
\begin{align*}
Qx + c + C^\top\lambda &= 0,\\
Cx - d &= 0.
\end{align*}
Equivalently, we obtain the linear KKT system
\begin{equation*}
\begin{pmatrix}
6 & 0 & 2 & 1 & 0\\
0 & 4 & 0 & 1 & 1\\
2 & 0 & 6 & 2 & 1\\
1 & 1 & 2 & 0 & 0\\
0 & 1 & 1 & 0 & 0
\end{pmatrix}
\begin{pmatrix}
x_1\\ x_2\\ x_3\\ \lambda_1\\ \lambda_2
\end{pmatrix}
=
\begin{pmatrix}
3\\ 2\\ -7\\ 0\\ 2
\end{pmatrix},
\end{equation*}
where the right-hand side comes from $Qx + C^\top\lambda = -c$ and $Cx = d$.

We now solve this linear system using Gaussian elimination.
The augmented matrix is
\begin{equation*}
\left[
\begin{array}{ccccc|c}
6 & 0 & 2 & 1 & 0 & 3\\
0 & 4 & 0 & 1 & 1 & 2\\
2 & 0 & 6 & 2 & 1 & -7\\
1 & 1 & 2 & 0 & 0 & 0\\
0 & 1 & 1 & 0 & 0 & 2
\end{array}
\right].
\end{equation*}
Applying Gaussian elimination (row operations), we reduce it to row-echelon form and then to reduced row-echelon form.
The resulting matrix is
\begin{equation*}
\left[
\begin{array}{ccccc|c}
1 & 0 & 0 & 0 & 0 & -1\\
0 & 1 & 0 & 0 & 0 & 3\\
0 & 0 & 1 & 0 & 0 & -1\\
0 & 0 & 0 & 1 & 0 & 11\\
0 & 0 & 0 & 0 & 1 & -21
\end{array}
\right],
\end{equation*}
which is the reduced row-echelon form of the KKT system.

From this we read off the unique solution
\begin{equation*}
x_1^\star = -1,\qquad x_2^\star = 3,\qquad x_3^\star = -1,
\end{equation*}
and the corresponding Lagrange multipliers
\begin{equation*}
\lambda_1^\star = 11,\qquad \lambda_2^\star = -21.
\end{equation*}
One can verify that the constraints are satisfied:
\begin{align*}
x_1^\star + x_2^\star + 2x_3^\star &= -1 + 3 + 2(-1) = 0,\\
x_2^\star + x_3^\star &= 3 + (-1) = 2.
\end{align*}
Since $Q$ is positive definite, the quadratic objective is strictly convex and the KKT conditions are sufficient for optimality. Therefore $x^\star = (-1,3,-1)^\top$ is the unique minimizer of the problem.

The minimal objective value is
\begin{equation*}
f(x^\star)
= (-3\ -2\ 7)x^\star + \tfrac12 (x^\star)^\top Q x^\star
= 16.
\end{equation*}

\end{document}